\subsection{RedPajama、RefinedWeb、FineWeb：网页数据进入精加工时代}

LLaMA 的公开配方启发了 RedPajama 和 SlimPajama（Cerebras 对 RedPajama 做 MinHashLSH 去重后得到的 627B token 子集）。随后 RefinedWeb 和 FineWeb 又把网页数据处理进一步推进。

\textbf{RefinedWeb} 的核心观点是”web data is all you need”——只要清洗得够好，纯网页数据就能媲美多源混合。它使用 trafilatura 从 WARC 提取文本（而非依赖 WET），应用 Gopher 的规则过滤（有意避免 ML 过滤器以减少偏见），并用 MinHash 在 5-gram 上做模糊去重。最终发布了 600B tokens（总量 5T）。

\textbf{FineWeb} 最初是对 RefinedWeb 的复现尝试，但在过程中做了多项改进：处理了 95 个 Common Crawl 快照，加入了 URL 过滤、语言识别（英语概率 $> 0.65$）、Gopher + C4 规则过滤、MinHash 去重，以及邮箱和公共 IP 地址的匿名化（PII 保护）。最终产出 15T tokens。

\begin{importantbox}{网页数据工程的底层变化}
早期重点是”尽量多抓”。后来重点变成”怎样抓得更像人想读的文本，而且不把隐私和噪声一起放大”。关键转折是：当 RefinedWeb 证明精心清洗的纯网页数据可以和多源混合相竞争时，整个社区开始认真对待数据清洗本身的工程价值。
\end{importantbox}

